\chapter{Camino hacia un circuito integrado}
\label{ch:ch6_asic}

% [BORRADOR v1] Material de partida: docs/sintesis_resultados.md, docs/via_asic.md,
% synth/sintetiza.py, results/sintesis_sg13g2.csv.

El motor está diseñado para acabar siendo un circuito integrado dedicado que
sirve claves a un microcontrolador, y la FPGA es su vehículo de
verificación, no su destino. Este capítulo toma el mismo VHDL que pasa los
bancos de pruebas del Capítulo~\ref{ch:ch5_verification}, lo sintetiza
sobre una biblioteca de celdas estándar abierta y expone cuánto costaría una
tirada de fabricación y qué tendría que cambiar en la fuente de entropía. Un
envío a fabricación (tape-out) queda explícitamente fuera del alcance de este TFM
--- no por coste ni por área, como se verá, sino por plazos --- y todo lo que
aquí aparece se etiqueta como lo que es: una cota inferior del área obtenida de
la síntesis lógica, con los factores de crecimiento hasta el área del núcleo
tomados de datos publicados.

\section{Flujo de síntesis con celdas estándar}
\label{sec:synthflow}

El flujo es reproducible a propósito con herramientas libres. Yosys lee el
VHDL a través del front end de GHDL, sintetiza cada bloque aplanado, mapea
los flip-flops y la lógica sobre el fichero liberty de la biblioteca objetivo
e informa del número de celdas y del área. La biblioteca es el \ac{PDK}
abierto de IHP, SG13G2, un proceso BiCMOS de 130~nm con una biblioteca de
celdas estándar mantenida, macros de SRAM compiladas, celdas de E/S y soporte
de OpenROAD. Se ha elegido porque es el único \ac{PDK} abierto que está vivo,
es europeo y es accesible a través de la misma afiliación a Europractice que
la universidad ya tiene. La puerta equivalente se lee del propio fichero
liberty y no de memoria:
$1~\mathrm{GE} = $ área de \texttt{sg13g2\_nand2\_1} $= 7{,}2576~\mu\mathrm{m}^2$.
No se aplica ninguna restricción de reloj: Yosys no la usa para el área, así
que las cifras no dependen de una frecuencia objetivo. Los osciladores en
anillo no se sintetizan; son el único bloque diseñado a nivel de transistor
(Capítulo~\ref{ch:ch3_entropy}).

\section{Área por bloque en puertas equivalentes}
\label{sec:area}

La Tabla~\ref{tab:synth} da el resultado de cada bloque. El motor digital
--- todo salvo el búfer de captura --- ocupa unos \textbf{46~kGE}.
Sintetizados por separado, el autenticador y el generador llevan cada uno
su propio núcleo de cifrado de una ronda por ciclo, de 16~kGE, y así
están medidas sus filas de la tabla; en el motor completo un solo núcleo
de 32 bits sirve a los dos, porque el secuenciador nunca los deja
funcionar a la vez (Sección~\ref{sec:opt}), de modo que las filas no
suman la cifra del motor. Aproximadamente: el AES compartido de 32 bits
6,7~kGE, el autenticador sin su núcleo 7, el generador sin él 22, los
tests de salud 1,3, muestreador, cruce y resets 0,4, \ac{I2C} 1,1, y
unos 7~kGE de secuenciador, registros y reparto del núcleo en el nivel
superior. El
búfer de captura se incluye para que la tabla esté completa, pero no debe
leerse como parte del producto: sus 46,8~kGE son una memoria de 4096 bits
mapeada sobre flip-flops porque el flujo no tiene macros de SRAM. En la FPGA
ocupa una block RAM sin coste de lógica; en silicio sería una SRAM compilada
o, con más sentido, un búfer mucho más pequeño, porque es una ruta de prueba
que se desactiva en producción.

\begin{table}[htbp]
\centering
\caption{Síntesis lógica sobre IHP SG13G2, aplanada, sin rutado.
1~GE $=7{,}2576~\mu\mathrm{m}^2$. Cota inferior del área \simu{}.}
\label{tab:synth}
\begin{tabular}{@{}lrrrr@{}}
\toprule
Bloque & Celdas & Flip-flops & Área ($\mu\mathrm{m}^2$) & kGE \\
\midrule
Núcleo AES-128, una ronda por ciclo (20 S-boxes, entrada de borrado) & 10\,571 & 262 & 116\,550 & 16,06 \\
Núcleo AES-128, ruta de 32 bits (4 S-boxes, 51 ciclos/bloque) & 3\,813 & 266 & 48\,434 & 6,67 \\
AES-CMAC (con su AES de una ronda) & 13\,289 & 912   & 169\,382 & 23,34 \\
CTR\_DRBG con df (con su AES de una ronda) & 19\,895 & 2\,357 & 277\,293 & 38,21 \\
Tests de salud (RCT, APT)      & 588     & 86    & 9\,137   & 1,26 \\
Muestreador ERO y divisor de 24 bits & 127 & 28  & 2\,371   & 0,33 \\
Cruce de dominio de reloj del bit & 11   & 7     & 408      & 0,06 \\
Sincronizador de reset (uno por dominio) & 3 & 3 & 147      & 0,02 \\
Esclavo I2C (con guarda de bus bloqueado) & 435 & 89 & 7\,729 & 1,07 \\
\midrule
Búfer de captura (4096 bits en flip-flops) & 12\,526 & 4\,140 & 339\,386 & 46,76 \\
Motor completo, anillos como cables, un solo AES de una ronda & 39\,879 & 8\,175 & 739\,061 & 101,83 \\
\textbf{Motor completo, un solo AES de 32 bits} & 32\,954 & 8\,179 & 671\,576 & \textbf{92,53} \\
\bottomrule
\end{tabular}
\end{table}

Los números son coherentes entre sí, y esa coherencia es la que permitió
detectar un error en el flujo: antes de aplanar, el informe devolvía para el
autenticador y para el generador solo el área de su submódulo AES. Un
autenticador que pese exactamente lo mismo que un AES desnudo es imposible,
así que el que fallaba era el informe, no el diseño. Tras aplanar, el
autenticador añade al AES 650 flip-flops (subclaves, estado de
encadenamiento, clave) y 7,6~kGE, y el generador añade 2095 flip-flops y
22~kGE, de los que cerca de la mitad son sus registros. La tabla es la
síntesis del motor con la versión~3 del mapa de registros, que añadió la
orden de borrado total, un contador de peticiones que satura en vez de
dar la vuelta, las reglas de exclusión de la captura y los
sincronizadores de reset; el generador creció 1,0~kGE por el contador
saturante y la ruta de borrado, y el motor completo 0,8~kGE, hasta
115,7~kGE con dos núcleos; compartir el núcleo lo dejó en 101,8, y el
núcleo de 32 bits en los 92,5~kGE de la tabla. El flujo de
esta tanda leyó el \ac{VHDL} a través de la propia salida Verilog de
GHDL y de la versión de Yosys instalable con \texttt{pip}, en lugar del
conector de GHDL; en el núcleo AES, que no cambió, los dos caminos
coinciden dentro del 0,9\,\%.

\section{Optimizaciones y sus ahorros esperados}
\label{sec:opt}

La síntesis es ingenua a propósito --- el \ac{RTL} se escribió para ser
verificable, no pequeño --- y apunta a tres optimizaciones, por orden de
rentabilidad del esfuerzo:

\begin{enumerate}
\item \textbf{Un único núcleo AES compartido --- hecho.} El secuenciador
del nivel superior ya serializa las órdenes: el generador y el
autenticador nunca funcionan a la vez. El núcleo vive ahora en el nivel
superior y se entrega al bloque que levanta su pulso de arranque; una
propiedad del banco del sistema comprueba que los dos nunca lo hacen en
el mismo ciclo. Ahorro medido: \textbf{14,1~kGE} (de 115,7 a 101,6), algo
por debajo de los 15,8 estimados, porque el reparto cuesta cerca de medio
kGE y la entrada de borrado que la revisión exigió al núcleo y al
generador cuesta otro. Un genérico conserva la disposición de dos
núcleos.
\item \textbf{Ruta de datos compacta para la S-box --- hecho.} Veinte
tablas S-box independientes se sintetizan en 16~kGE para el núcleo de
una ronda por ciclo; los diseños optimizados en área de la literatura van
desde 2,4~kGE para una ruta de datos de 8 bits a 226 ciclos por
bloque~\cite{moradi2011pushing} hasta unos 5~kGE para rutas de datos de
32 bits. Aquí la velocidad es irrelevante, así que se escribió un segundo
núcleo con ruta de 32 bits: cuatro S-boxes compartidas entre SubBytes y
el SubWord de la expansión de clave, un solo MixColumns, ShiftRows
aplicado como permutación de cableado del registro de estado al empezar
cada ronda, una columna por ciclo, 51 ciclos por bloque en vez de 11.
Pasa los mismos vectores oficiales, la cadena de mil bloques y la prueba
de zeroización, y el banco exige que su salida coincida con la del núcleo
de una ronda en cada bloque. Medido: \textbf{6,67~kGE frente a 16,06} en
el núcleo (ahorro de 9,4, frente a los 10 estimados; la diferencia son
los multiplexores de columna y el registro de fase) y \textbf{92,5 frente
a 101,8~kGE} en el motor completo, en la misma tanda de síntesis. El
núcleo de una ronda sigue disponible mediante un genérico.
\item \textbf{Menos registros en el generador.} De sus 2355 flip-flops, la
entropía capturada y el nonce (384), el material de semilla (256), la salida
(256) y los temporales de la función de derivación pueden reducirse
consumiendo la semilla como un flujo y compartiendo temporales; unos 800
flip-flops, cerca de 4~kGE.
\end{enumerate}

Las dos optimizaciones medidas llevaron el motor digital de 69~kGE (dos
núcleos de una ronda) a 46; la tercera lo dejaría en
\textbf{$\approx$42~kGE}. Es más que los 13 a 25~kGE que anticipaba el
estado del arte para motores de este tipo, y la tabla dice por qué: el
generador sin su núcleo son 22~kGE, la mitad registros, y el
secuenciador y los registros del nivel superior otros 7; llegar al rango
publicado exigiría repensar el almacenamiento del generador y no su
cifrador, lo que este trabajo deja planteado y no hecho. Con el factor
de crecimiento publicado de $\approx$2,35 del área de celdas al área del
núcleo para un AES en esta misma biblioteca, 42~kGE son unos
0,72~mm$^2$ de núcleo, que siguen cabiendo varias veces en el bloque
académico más pequeño (Sección~\ref{sec:shuttle}). Los dos primeros
ahorros son medidas; el tercero es una estimación y se recoge como
trabajo futuro, no se da por hecho.

\section{Lo que necesita el oscilador en anillo en silicio}
\label{sec:ro_silicon}

Los bloques digitales se llevan a silicio con el flujo estándar. La fuente
de entropía no, y la literatura que ha dado ese salto --- en particular un
estudio de 2024 que caracteriza el mismo modelo de anillo en una Artix-7 y
en un circuito FD-SOI de 28~nm~\cite{benea_2024_flicker} --- plantea tres
puntos que el diseño del chip debe asumir.

En primer lugar, el anillo se reconstruye con celdas inversoras en lugar de
tablas de consulta, y su jitter cambia en valor y en naturaleza: el ruido
de parpadeo gana peso en los nodos más recientes, y el estudio encuentra, en
contra de lo esperado, que contarlo como fuente legítima puede aumentar el
tasa de bits para la misma entropía. El modelo estocástico y el estimador del
Capítulo~\ref{ch:ch3_entropy} se trasladan sin cambios; la $\sigma$ medida
no. En segundo lugar, la variación de proceso hace que cada dado sea
distinto, así que la frecuencia y el jitter de cada anillo deben calibrarse
chip a chip; los diseños totalmente digitales publicados que superan
barridos de tensión y temperatura lo hacen con un anillo configurable y un
lazo de ajuste~\cite{yang_2016_jssc}. En tercer lugar, y lo más importante
para un dispositivo de seguridad, inyectar una frecuencia en la alimentación
engancha los anillos y elimina el jitter: una señal de 500~mV en el raíl
redujo el espacio de claves de un microcontrolador seguro comercial de
$2^{64}$ a unas 3300~\cite{markettos_2009_frequency}, y unos pocos cientos de
microvatios de campo radiado destruyeron la estadística de un generador de
cincuenta anillos~\cite{bayon_2012_contactless}. El mismo lazo de
calibración, junto con el filtrado de la alimentación, es la defensa
publicada, y los tests de salud de la Sección~\ref{sec:health} son el
detector de último recurso. Nada de esto está en el prototipo en FPGA; es el
trabajo que añade una implementación en silicio.

\section{Coste y plazos de una lanzadera académica}
\label{sec:shuttle}

La fabricación es más accesible de lo que sugiere el peso de la palabra. La
Universidad de Sevilla y el IMSE-CNM son miembros de Europractice, cuya
lista de precios académicos es pública: un bloque mini@sic en UMC~180~nm
cuesta 3\,430~\texteuro{}, uno en TSMC~65~nm 3\,691~\texteuro{}, e IHP SG13G2
se ofrece a 5\,110~\texteuro{}/mm$^2$ con una tirada multiproyecto gratuita
para diseños de código abierto de menos de 2~mm$^2$~\cite{europractice_prices2026}.
Un motor de 0,3--0,4~mm$^2$ cabe muchas veces en el bloque más pequeño de
cualquiera de ellos; el chip estaría limitado por sus pads, no por su
lógica. El propio grupo de seguridad hardware del instituto ya ha llevado
primitivas de oscilador en anillo desde FPGAs Xilinx~7-series hasta silicio
de 65~nm caracterizado, lo que hace que continuar este trabajo sea cuestión
de sumarse a una línea existente más que de abrir una.

Lo que no cabe en un trabajo fin de máster son los plazos: desde el envío a
fabricación (tape-out) hasta tener muestras encapsuladas pasan de meses a
más de un año, el precio no incluye el encapsulado y la caracterización de
los dados devueltos es un proyecto en sí misma. Como comparación, un
elemento seguro comercial como el ATECC608B cuesta 0,90~\$ por unidad suelta
y se vende con una validación NIST de la fuente de entropía~\cite{microchip_atecc608b_ds};
dos de sus competidores directos no mencionan SP~800-90 en ninguna parte de
sus hojas de datos. Por tanto, la posición honesta de este trabajo no es
competir en precio sino en \emph{transparencia de la entropía}: cada bit de
una clave producida por este motor puede trazarse hasta un jitter medido y
una cota declarada, que es precisamente lo que una pieza de caja negra no
puede ofrecer.
